\section{Conclusion}

We studied whether the factorization gap in discrete diffusion decoding can be addressed selectively rather than globally. On a controlled synthetic scaffold, Adaptive Dependency Routing Diffusion improves over fully factorized decoding while using non-factorized handling on only a minority of tokens. The gains are especially clear on dependency-sensitive and longer-context settings. However, the best current same-budget surrogate is not the full adaptive policy but a simpler coupled-only alternative. That negative result matters: it shows that selective dependency allocation is worth studying, but also that the controller itself still needs stronger justification.

The most defensible conclusion is therefore modest. Selective dependency handling appears feasible and measurable. A budgeted routing view of diffusion decoding is plausible. But the present evidence is still a pilot study, and the next serious test must happen on real diffusion language models.
